% Standalone document for the thesis progress/roadmap section.
% This content is NOT part of the thesis (main.tex); it tracks completed and
% remaining work and is compiled on its own via `make roadmap`.
\documentclass[11pt]{article}

\usepackage[margin=1in]{geometry}
\usepackage[colorlinks]{hyperref}
\usepackage{amsmath,amssymb}
\usepackage{cite}

\title{Thesis Progress: Completed and Remaining Work}
\author{Zhiyang Chen}
\date{\today}

\begin{document}
\maketitle

\section{Completed Work and Remaining Work}

% This thesis studies two families of high-stakes software systems operating in untrusted environments.
% The first family is on-chain: deployed smart contracts that execute directly on a blockchain and must tolerate arbitrary user interactions in public ecosystems.
% The second family is off-chain: AI-assisted coding systems that operate over external tools and online information, where behavior may be influenced by opaque training data and weakly trusted content.
% Within this framing, the on-chain track studies runtime protection for deployed smart contracts,
% while the off-chain track studies security risks introduced through AI-assisted software development.
% Within this scope, the dissertation is grounded in three completed projects and one planned project.
% The completed projects establish the core technical results; the planned project defines the final extension.
% Conceptually, the section below is organized as a status-oriented roadmap: what has been completed,
% what evidence those completed projects already provide, and what concrete work remains to close the thesis.


\subsection{Overall Plan and Current Status}

The completed and ongoing work follows the timeline below.
\begin{itemize}
	\item Oct 2022 -- Oct 2023: \textsc{Trace2Inv} research and implementation;
	later revised and published at FSE 2024.
	\item Nov 2023 -- Sept 2024: \textsc{CrossGuard} research and implementation;
	published at ICSE 2026.
	\item Apr 2025 -- Oct 2025: \textsc{Scam2Prompt} research and large-scale evaluation;
	submitted to ICML 2026.
	\item Current remaining phase: one final project on AI-era scam and malware 
	distribution in developer pipelines.
\end{itemize}

At this point, I have already finished work on
(i) invariant-centered runtime defense,
(ii) control-flow-integrity runtime defense, and
(iii) auditing of malicious artifacts in LLM-assisted coding.
The remaining work is to study how poisoned external content 
retrieved by coding agents 
can propagate into harmful coding outcomes, and how to 
defend against this pipeline.


\subsection{Completed Work}

\noindent \textbf{Trace2Inv (completed; published at FSE 2024).}
\textsc{Trace2Inv}~\cite{chen2024demystifying} was developed from Oct 2022 to Oct 2023.
The project addresses a central gap in contract protection: manually specified invariants, namely rules that should remain true during normal execution, are often incomplete,
and static specifications frequently miss protocol-level behavioral assumptions observed only in deployment.
\textsc{Trace2Inv} learns invariants from decoded transaction history and enforces them as runtime guards, that is, checks inserted to stop suspicious executions while the contract is running.
Its design supports 23 invariant templates across 8 categories, including access control, time locks,
gas usage, re-entrancy, oracle slippage, special storage, money flow, and data flow.
The pipeline contains four stages: invocation-tree construction, invariant-relevant data extraction,
template-based invariant generation, and runtime enforcement.
In evaluation, the resulting guard policies achieved strong protection with low false-positive rates (0.32\%).
This project establishes the thesis's first technical pillar: execution-trace-grounded runtime invariants
as a practical security primitive for deployed protocols.

\noindent \textbf{CrossGuard (completed; published at ICSE 2026).}
\textsc{CrossGuard}~\cite{chen2025secure} was developed from Nov 2023 to Sept 2024.
This project studies an orthogonal runtime-defense question: whether protocol attacks can be blocked
through control-flow integrity, meaning that functions are allowed to execute only in legitimate call patterns, rather than vulnerability-specific signatures.
\textsc{CrossGuard} enforces function-level control-flow integrity through pre-deployment instrumentation
and runtime invocation validation.
The system uses whitelist-based policies derived from benign invocation structure,
including independent invocation validation, read-only invocation filtering,
runtime detection of read-only functions, and restore-on-exit storage handling.
A key engineering objective is deployability: guard enforcement is performed while preserving compatibility
with existing DeFi contracts, namely blockchain programs for services such as trading and lending, and while limiting overhead.
Empirically, \textsc{CrossGuard} achieved strong protection with 15.5\% average gas overhead, that is, additional transaction cost paid on the blockchain.
Relative to \textsc{Trace2Inv}, this project advances the thesis from property-level constraints
to interaction-structure integrity, broadening defense coverage under realistic deployment constraints.

\noindent \textbf{Scam2Prompt (completed; accepted to ICML 2026).}
\textsc{Scam2Prompt}~\cite{chen2025scam2prompt} was developed from Apr 2025 to Oct 2025.
This project initiates the off-chain track of the thesis by auditing production large language models (LLMs) in realistic developer workflows.
Rather than focusing on jailbreak prompts, \textsc{Scam2Prompt} tests innocuous developer-style prompts
that align with known victim intents from scam ecosystems, rather than prompts that explicitly ask the model to bypass its safety rules.
Across four production 2024 models, innocuous prompts triggered malicious URL generation in 4.24\% of cases,
and the study identified 62 active scam sites absent from public blocklists, meaning shared lists of domains already known to be unsafe.
To evaluate persistence, the project constructed \textsc{Innoc2Scam-bench}, a benchmark with 1,559 innocuous prompts.
Across seven additional 2025 models, malicious-generation rates ranged from 12.7\% to 43.8\%,
while contemporary guardrails detected less than 0.3\% of these cases.
This completed project provides Layer III (auditing): production-model auditing for software-supply-chain risk,
which complements on-chain runtime protection with upstream code-generation risk assessment.

Taken together, these completed projects establish the first three layers of the thesis:
Layer I (invariants), Layer II (control-flow integrity), and Layer III (auditing).
Layer IV is the planned remaining work on search- and documentation-level poisoning of coding agents.




\subsection{Work Remaining}

\noindent \textbf{Position in the thesis.}
This remaining work constitutes the fourth layer of the dissertation.
Layers I and II study runtime defenses for smart contracts, including invariant enforcement 
and control-flow integrity.
Layer III studies harmful artifacts in code suggested by LLMs.
Layer IV moves one step further upstream by examining an external attack surface:
search- and documentation-level poisoning of coding agents.

\paragraph{Motivation.}
Modern coding agents increasingly rely on external resources such as official documentation,
API references, tutorials, and search results, when solving software engineering tasks.
While this substantially improves utility, it also introduces a new attack surface:
an adversary may influence agent behavior by manipulating the external 
content that the agent retrieves and trusts,
without changing model weights or injecting harmful prompts.

The core question is therefore end-to-end:
when poisoned content is present in retrieval channels, how often is it
(1) retrieved,
(2) treated as credible,
and (3) translated into concrete harmful actions?
This question is especially important because human developers do not 
treat all external sources equally.
Experienced developers often distinguish between mature and trustworthy libraries,
less reliable packages, unofficial tutorials, and low-quality discussions,
and they adjust their reliance accordingly. 
Experienced human developers might choose to ignore or discount 
less reliable sources, and implement their own checks or artifacts despite 
the extra effort and engineering time that doing so requires. 

However, it remains unclear whether coding agents exhibit a similar notion of source reliability,
or whether poisoned search results or documentation can alter their preferences 
and induce unsafe decisions.
This work therefore studies the interaction between an agent's prior tendencies and 
the external evidence it retrieves,
with a focus on when that interaction becomes security-critical.

\paragraph{Problem statement.}
This work studies \emph{external-content poisoning against coding agents}.
The attacker controls a subset of externally visible content,
such as documentation pages, API references, tutorial fragments, issue discussions, 
package descriptions, or search snippets,
while the victim is a coding agent that retrieves such content during task solving.
The attacker's goal is to induce harmful outcomes, including insecure patches,
plausible but incorrect patches, unsafe dependency choices, or unsafe recommendations in the final output.
The central research challenge is to measure and explain compromise across the full decision pipeline:
\emph{retrieval}, \emph{trust}, and \emph{action}.

\paragraph{Research questions.}
\begin{itemize}
    \item \textbf{RQ1.}
    To what extent are open-source coding agents (e.g., OpenCode) vulnerable to poisoned search results
    and documentation pages during realistic software-engineering tasks that require external lookup?

    \item \textbf{RQ2.}
    How does compromise vary with attacker budget, retrieval rank, corroborating poisoned sources,
    source reputation, and detectability constraints?
    In particular, do coding agents distinguish between sources of different reliability,
    or can poisoning override such preferences?

    \item \textbf{RQ3.}
    Which practical and deployable defenses reduce compromise
    while preserving benign task performance and agent utility?
\end{itemize}

\paragraph{Threat model.}
We assume a black-box attacker, meaning an attacker who cannot directly 
modify the internal model or agent implementation.
More specifically, the attacker cannot modify model weights, namely the learned internal parameters of the model,
agent system prompts, namely the hidden instructions that guide the agent's behavior, or the agent's internal implementation.
However, the attacker can control part of the external content visible through search and documentation channels
by publishing or manipulating publicly visible web content.
The victim is a coding agent that is allowed to browse or retrieve such content while solving a task.
The primary outcomes of interest are:
(1) security-degrading patches,
(2) unsafe dependency selections or API usage,
and
(3) unsafe recommendations or citations in final outputs.

\paragraph{Experimental design.}
The empirical study will evaluate at least two open-source coding-agent stacks
(e.g., an OpenCode-style system and a SWE-agent-style system)
on task suites that genuinely require external lookup.
To ensure reproducibility, the main benchmark will be executed over a controlled local retrieval corpus
that combines benign resources with attacker-controlled content, where a retrieval corpus is the document collection from which the agent searches and retrieves information.
This design enables systematic manipulation of rank, visibility, source agreement, and source reputation.

The evaluation will consider several poisoning strategies, including:
(1) \emph{rank poisoning}, where attacker-controlled content is made highly visible;
(2) \emph{content poisoning}, where retrieved documents contain misleading or adversarial technical guidance; and
(3) \emph{consensus poisoning}, where multiple poisoned sources appear to corroborate one another.
Measurements will be collected at each stage of the pipeline,
including whether poisoned content is retrieved, whether it is cited or followed,
and whether it leads to harmful code actions.
The study will also measure the tradeoff between attack resistance and normal task success.

\paragraph{Defenses and expected deliverables.}
The defense study will compare lightweight and deployable mechanisms, including:
source-level filtering,
instruction-level separation,
cross-source agreement gating,
and post-generation security checking.
The expected deliverables include:
\begin{itemize}
    \item a benchmark for search- and documentation-level 
    poisoning of coding agents,
    \item deployment-oriented guidance for reducing upstream 
    risk in AI-assisted software systems.
\end{itemize}

\paragraph{Scope.}
To keep this dissertation layer focused and publishable,
the work targets coding agents and documentation/search poisoning only.
It does not attempt to cover all browser-agent environments,
all prompt-injection settings,
or all forms of model poisoning.
This scope aligns Layer IV with the dissertation's 
overall goal of developing practical 
techniques for hardening the software-development 
pipeline against security threats.



% \section{Organization}

% The thesis is organized as follows.
% Chapter~\ref{sec:preli} defines the formal preliminaries and 
% unified notation, including Section~\ref{sec:background} on execution-semantic background.

% Across these chapters, we follow a consistent structure: problem context,
% formal model, methodological choices, and implications.
% This organization is designed to keep cross-layer reasoning explicit and technically auditable.


\bibliographystyle{abbrv}
\bibliography{main}

\end{document}
